\chapter{Power: Where the Energy Goes}
\label{ch:power}

Power limits modern chips more than anything else. A GPU could be faster if it could burn more power, but it would melt (Chapter~\ref{ch:heat}), and a phone could be faster but its battery would die. This chapter explains exactly where a chip's power goes, derives the formula behind the assessment question ``if frequency and voltage double, what happens to dynamic power?'', and surveys the techniques engineers use to save power.

\section{Three kinds of power}

\begin{enumerate}
  \item \term{Dynamic power}: spent charging and discharging capacitance every time a signal switches. Usually the largest part when a chip is busy.
  \item \term{Short-circuit power}: during each transition, both the pull-up and pull-down are briefly on together, and current flows straight from the supply to ground.
  \item \term{Leakage power} (static power): transistors that are ``off'' still leak a little current, all the time, even when nothing switches.
\end{enumerate}

\section{Dynamic power}

Consider an inverter driving a load capacitance $C$ (Figure~\ref{fig:dynpower}).

\begin{figure}[h]
\centering
\begin{circuitikz}
  \draw (0,0) node[pmos](p){} (0,-2.2) node[nmos](n){};
  \draw (p.S) node[vdd]{$\VDD$};
  \draw (n.S) node[ground]{};
  \draw (p.D) -- (n.D);
  \draw (p.G) -- (n.G);
  \draw ($(p.G)!0.5!(n.G)$) to[short,-o] ++(-0.8,0) node[left]{In};
  \draw ($(p.D)!0.5!(n.D)$) to[short,*-] ++(1.6,0) coordinate(o) to[C, l=$C$] ++(0,-1.6) node[ground]{};
  \draw[->,red!70!black,thick] (0.45,0.2) .. controls (1.2,0.2) .. (1.55,-0.8);
  \node[red!70!black,font=\small,right] at (1.5,0.3) {charge: $CV^2$ from supply};
  \draw[->,blue!70!black,thick] (1.55,-1.5) .. controls (1.0,-2.6) .. (0.45,-2.4);
  \node[blue!70!black,font=\small,right] at (1.6,-2.6) {discharge to ground};
\end{circuitikz}
\caption{Each 0-to-1 output transition draws charge from the supply; each 1-to-0 transition dumps it to ground.}
\label{fig:dynpower}
\end{figure}

\paragraph{Charging (output 0 $\to$ 1).} The PMOS connects the capacitor to $\VDD$. Charge $Q = C\VDD$ flows out of the supply at voltage $\VDD$, so the supply delivers energy $Q\VDD = C\VDD^2$. Of that, $\tfrac12 C\VDD^2$ ends up stored in the capacitor (Chapter~\ref{ch:capind}); the other half is turned into heat in the PMOS's resistance.

\paragraph{Discharging (output 1 $\to$ 0).} The NMOS connects the capacitor to ground, and the stored $\tfrac12 C\VDD^2$ is turned into heat in the NMOS.

So every full cycle of charge and discharge consumes $C\VDD^2$ from the supply. If the node does this $f$ times per second, the power is $C\VDD^2 f$. Most nodes do not switch every cycle, so we multiply by an \term{activity factor} $\alpha$: the fraction of clock cycles in which the node makes a 0-to-1 transition.

\begin{center}
\fbox{\parbox{0.6\linewidth}{\centering\vspace{0.2em}
$\displaystyle P_{dynamic} = \alpha\, C\, \VDD^2\, f$\vspace{0.4em}}}
\end{center}

\begin{itemize}
  \item $\alpha$: activity factor. A clock rises once every cycle, so $\alpha = 1$. Typical data nodes have $\alpha$ around 0.1 to 0.25.
  \item $C$: the total capacitance that switches.
  \item $\VDD$: supply voltage, and it is \emph{squared}.
  \item $f$: clock frequency.
\end{itemize}

\begin{bigidea}
Dynamic power is proportional to the \emph{square} of the voltage and only linearly to frequency and capacitance. Lowering the voltage is the most powerful way to save power.
\end{bigidea}

\begin{worked}{a chip's dynamic power}
A block has 1\,nF of total switching capacitance, an average activity factor of 0.2, a 1.0\,V supply and a 1\,GHz clock.
\[
  P = 0.2 \times 10^{-9} \times 1.0^2 \times 10^9 = 0.2~\text{W}.
\]
\end{worked}

\section{Scaling questions}
\label{sec:powerscaling}

Assessment questions change one or more of the variables and ask for the new power. Since $P \propto \alpha C V^2 f$, multiply the factors, remembering to square the voltage factor.

\begin{center}
\begin{tabular}{lll}
\toprule
Change & New dynamic power & Working \\
\midrule
$f \times 2$ & $\times 2$ & linear in $f$ \\
$V \times 2$ & $\times 4$ & $2^2$ \\
$f \times 2$ \textbf{and} $V \times 2$ & $\times 8$ & $2 \times 2^2$ \\
$V \times 0.5$ & $\times 0.25$ & $0.5^2$ \\
$V$ up 10\%, $f$ up 20\% & $\times 1.452$ (+45\%) & $1.1^2 \times 1.2$ \\
$(f, 2V) \to (4f, 3V)$ & $\times 9$ & $4 \times (3/2)^2$ \\
$V \times 0.9$, $f \times 0.9$ & $\times 0.729$ & $0.9^2 \times 0.9$ \\
$C \times 2$ (twice the logic) & $\times 2$ & linear in $C$ \\
\bottomrule
\end{tabular}
\end{center}

\begin{onoa}
``If the clock frequency doubles and the supply voltage doubles, what is the dynamic power?'' was an Ignite assessment question. Answer: \textbf{8 times} the original ($2$ for frequency times $2^2 = 4$ for voltage). Reported variants include voltage up 10\% with frequency up 20\% ($\times 1.45$) and a change from $(f, 2V)$ to $(4f, 3V)$ ($\times 9$; note that the voltage factor is $3V/2V = 1.5$, squared to 2.25).
\end{onoa}

\begin{trap}
Watch for the starting point. In ``$(f, 2V) \to (4f, 3V)$'' the voltage does not triple; it goes from $2V$ to $3V$, a factor of 1.5. Always divide new by old for each variable before multiplying.
\end{trap}

\section{Power versus energy}

\term{Power} is energy per second. \term{Energy} is what drains a battery. The energy for one operation is
\[
  E_{op} = \alpha C \VDD^2,
\]
with no $f$ in it. So:
\begin{itemize}
  \item Running the same work at double the frequency (same voltage) \emph{doubles the power} but uses \emph{the same energy} per operation; the work just finishes sooner.
  \item Lowering the voltage reduces \emph{both} power and energy per operation.
\end{itemize}
This is why ``slow down to save energy'' only works if the voltage comes down too.

\section{Short-circuit power}

During an input transition, there is a moment when the input is between $\Vtn$ and $\VDD - |\Vtp|$ and both transistors are on. Current flows straight from $\VDD$ to ground. The slower the input transition, the longer this lasts, which is one reason designers keep signal transitions sharp. In a well-designed chip, short-circuit power is usually 10\% or less of dynamic power.

\section{Leakage power}
\label{sec:leakage}

Even with no switching at all, current trickles through transistors that are supposed to be off. Leakage power is $P_{leak} = I_{leak}\,\VDD$. The main sources:

\begin{itemize}
  \item \textbf{Subthreshold leakage}: an ``off'' transistor still conducts a little, and this current grows \emph{exponentially} as $\Vt$ decreases (Chapter~\ref{ch:mosfet}) and as temperature increases (Chapter~\ref{ch:semis}). It is usually the largest leakage source.
  \item \textbf{Gate leakage}: electrons tunnel through extremely thin gate oxide. High-$k$ gate materials were introduced largely to reduce it.
  \item \textbf{Junction leakage}: small reverse current through the source/drain diodes.
\end{itemize}

Two transistors that are both off and in series leak much less than one (the \term{stack effect}), because the node between them rises slightly and turns both off more firmly.

As transistors shrank, thresholds had to fall to keep them fast, and leakage grew from negligible to a large fraction of total power. Leakage also feeds on heat: a hotter chip leaks more and gets hotter still (Chapter~\ref{ch:heat}).

\section{Saving power}

\begin{center}
\begin{tabularx}{\linewidth}{>{\raggedright\arraybackslash}p{3.3cm}>{\raggedright\arraybackslash}p{2.3cm}>{\raggedright\arraybackslash}X}
\toprule
Technique & Reduces & How \\
\midrule
Clock gating & dynamic & Stop the clock to idle registers. The clock network has $\alpha = 1$ and is a big part of dynamic power, so this is the most common technique. \\
Power gating & leakage (and dynamic) & Cut off the supply to an idle block with a large ``sleep'' transistor (a header PMOS or footer NMOS). \\
DVFS & dynamic and leakage & Dynamic voltage and frequency scaling: lower $V$ and $f$ together when full speed is not needed. \\
Multi-$\Vt$ cells & leakage & High-$\Vt$ (low-leakage) transistors on paths that are not timing-critical. \\
Multiple voltage domains & dynamic & Run slower blocks at a lower $\VDD$; level shifters translate signals between domains. \\
Body biasing & leakage or delay & Reverse body bias raises $\Vt$ (less leakage); forward bias lowers it (faster). \\
Reducing activity & dynamic & Avoid useless toggling: operand isolation, Gray-coded counters, removing glitches. \\
Parallelism & dynamic & Two slower units at a lower voltage can beat one fast unit (below). \\
\bottomrule
\end{tabularx}
\end{center}

\paragraph{Clock gating done safely.} Simply ANDing a clock with an enable signal can produce glitches if the enable changes while the clock is high. The standard \term{integrated clock gating} (ICG) cell captures the enable in a latch that is transparent only while the clock is low, then ANDs the latched enable with the clock. The enable can only change while the clock is low, so the gated clock is always clean.

\paragraph{Power gating trade-offs.} A power-gated block loses its state, so registers that must survive need special \term{retention} flip-flops (or the state is saved elsewhere). Outputs of a powered-off block float, so \term{isolation cells} clamp them to safe values. Waking up takes time and causes a surge of current as the block's capacitance recharges.

\begin{worked}{parallelism saves power}
One processing unit runs at frequency $f$ and voltage $V$. Replace it with two identical units, each running at $f/2$ (same total throughput), which lets the voltage drop to $0.7V$. Ignoring overhead, the capacitance doubles:
\[
  \frac{P_{new}}{P_{old}} = \frac{2C \times (0.7V)^2 \times f/2}{C V^2 f} = 0.49.
\]
Same work, about half the power. This is a big part of why GPUs, with thousands of moderately clocked cores, are so energy efficient for parallel work.
\end{worked}

\begin{worked}{how much power is left after power gating?}
In an idle mode, a design power-gates most of its logic. The numbers (with assumptions stated):
\begin{itemize}
  \item Standard cells: 40\,W dynamic and 10\,W leakage when active. 10\% of the cells must stay on (always-on logic); the other 90\% are gated with 80\% efficiency (gating removes 80\% of their leakage).
  \item RAMs: 20\,W dynamic and 6\,W leakage. All are gated with 90\% efficiency.
  \item Gated blocks do no switching, so their dynamic power becomes zero; the always-on cells keep their share of dynamic power.
\end{itemize}
Standard-cell leakage left: $10 \times [0.1 + 0.9 \times (1 - 0.8)] = 10 \times 0.28 = 2.8$\,W.
RAM leakage left: $6 \times (1 - 0.9) = 0.6$\,W.
Dynamic left: $0.1 \times 40 = 4$\,W from the always-on cells; RAMs 0\,W.
Total in idle: $2.8 + 0.6 + 4 = 7.4$\,W, down from 76\,W.
\end{worked}

\begin{onoa}
A reported question gives power-gating efficiencies for standard cells (80\%) and RAMs (90\%), with 10\% of cells always on, and asks for the remaining dynamic and leakage power. The exact answer depends on the numbers given, but the method is the one above: split each block into always-on and gated parts, remove the stated fraction of leakage from the gated part, and remove the gated part's dynamic power.
\end{onoa}

\section{The bigger picture: Dennard scaling and dark silicon}

For decades, shrinking transistors brought a free lunch called \term{Dennard scaling}: smaller transistors could run at lower voltage, so power per unit area stayed constant while chips got faster. Around 2005 this broke down. Voltage could no longer drop much, because lower thresholds would make leakage explode. Power density started rising with each generation, clock frequencies stopped climbing (most processors have stayed around 3 to 5\,GHz since), and the industry turned to parallelism: multiple cores, and massively parallel GPUs.

Today a chip often cannot run all of its transistors at full speed at once without exceeding its power and cooling budget; some must be idle or slowed at any moment. This is called \term{dark silicon}. A chip's \term{thermal design power} (TDP) is the sustained power its cooling system is designed to remove.

\section{Summary}

\begin{itemize}
  \item Dynamic power $= \alpha C \VDD^2 f$: each charge-discharge cycle costs $C\VDD^2$; half is lost charging and half discharging.
  \item Scaling: multiply the factor changes, squaring the voltage factor. $f\times2$ and $V\times2$ gives $\times8$.
  \item Energy per operation $= \alpha C \VDD^2$, independent of $f$. Only lowering voltage saves energy per operation.
  \item Short-circuit power comes from both transistors being on during transitions; sharp edges minimize it.
  \item Leakage grows exponentially as $\Vt$ falls and temperature rises.
  \item Clock gating cuts dynamic power; power gating cuts leakage; DVFS cuts both.
  \item Parallel units at lower voltage can do the same work for less power.
  \item The end of Dennard scaling ended rising clock speeds and pushed the industry toward parallelism.
\end{itemize}

\section{Exercises}

\begin{enumerate}
  \item A block has $\alpha = 0.2$, $C = 1$\,nF, $\VDD = 1.0$\,V and $f = 1$\,GHz. Find its dynamic power.
  \item Clock frequency and supply voltage both double. By what factor does dynamic power change?
  \item Supply voltage rises 10\% and frequency rises 20\%. What is the percentage change in dynamic power?
  \item A chip moves from frequency $f$ at voltage $2V$ to frequency $4f$ at voltage $3V$. By what factor does dynamic power change?
  \item Voltage and frequency are both reduced by 10\%. What fraction of the original dynamic power remains?
  \item A processor runs the same program twice as fast at the same voltage. What happens to its power, and to the energy the program uses?
  \item Which kind of power does clock gating reduce? Which does power gating reduce?
  \item Does leakage power rise or fall with temperature?
  \item What is the activity factor of a clock net?
  \item Compare one core at frequency $f$ and 1.0\,V with two cores at $f/2$ and 0.6\,V each. What is the power ratio, ignoring overhead?
  \item Why can't a designer simply AND a clock with an enable signal to gate it?
\end{enumerate}

\subsection*{Answers}
\begin{enumerate}
  \item $0.2 \times 10^{-9} \times 1 \times 10^9 = 0.2$\,W.
  \item $2 \times 2^2 = 8$.
  \item $1.1^2 \times 1.2 = 1.452$, an increase of about 45\%.
  \item $4 \times (3/2)^2 = 9$.
  \item $0.9^2 \times 0.9 = 0.729$, about 73\%.
  \item Power doubles; energy for the program stays the same (energy per operation does not depend on frequency).
  \item Clock gating: dynamic power. Power gating: leakage power (and the gated block's dynamic power).
  \item Rise (exponentially).
  \item 1.
  \item $2 \times 0.6^2 \times 0.5 = 0.36$: about 36\% of the original power.
  \item If the enable changes while the clock is high, the AND output can produce a shortened pulse or glitch, which flip-flops would treat as a clock edge. A latch-based clock-gating cell allows the enable to change only while the clock is low.
\end{enumerate}
